\begin{name}
	{ÔN TẬP KIỂM TRA CUỐI HỌC KÌ 1}
	{TOÁN 10}
	{LỚP TOÁN THẦY PHÁT}
	{\thoigian}
\end{name}

%================================PHẦN I

%================================PHẦN I
\caulc
\Opensolutionfile{ans}[ans/0-HK1-CD-3-DongAnh-HaNoi-2324-Phan-I]

%[0-HK1-CD-3-DongAnh-HaNoi-2324]%[VN-MT-6, VN-22]

%%==========Câu 1
\begin{ex}%[0-HK1-CD-3-DongAnh-HaNoi-2324]%[VN-MT-6, VN-22]%[0D3N1-2]
	Tập xác định của hàm số $y=\sqrt{x-2}+\dfrac{2}{x+5}$ là
	\choice
	{$\mathscr{D}=\left(2;+\infty\right)\setminus\{5\}$}
	{$\mathscr{D}=\left[2;+\infty\right)$}
	{\True $\mathscr{D}=\left[2;+\infty\right)\setminus\{5\}$}
	{$\mathscr{D}=\left(-5;+\infty\right)$}
	\loigiai{
		Ta có 
		$\heva{&x-2\geq 0\\&x+5\ne0} \Leftrightarrow \heva{&x\geq2\\&x\ne 5.}$\\
		Tập xác định của hàm số là $\mathscr{D}=\left[2;+\infty\right)\setminus\{5\}$.
		
	}
\end{ex}

%%==========Câu 2
\begin{ex}%[0-HK1-CD-3-DongAnh-HaNoi-2324]%[VN-MT-6, VN-22]%[0D3H1-1]
	Parabol $y=a{x^2}+bx+2$ đi qua hai điểm $M(1;5)$ và $N(-2;8)$ có phương trình là
	\choice
	{$y=2x^2+2x+2$}
	{\True $y=2x^2+x+2$}
	{$y=x^2+2x+2$}
	{$y=x^2+x+2$}
	\loigiai{
		Thay tọa độ $M(1;5)$ vào phương trình Parabol ta có $5=a+b+2 \Leftrightarrow a+b=3.\qquad (1)$\\
		Thay tọa độ $M(-2;8)$ vào phương trình Parabol ta có $8=4a-2b+2 \Leftrightarrow 2a-b=3.\qquad (2)$\\
		Giải hệ phương trình $(1)$ và $(2)$ ta có $a=2$; $b=1$.\\
		Vậy Parabol đi qua hai điểm $M$ và $N$ có phương trình là $y=2x^2+x+2$.
	}
\end{ex}

%%==========Câu 3
\begin{ex}%[0-HK1-CD-3-DongAnh-HaNoi-2324]%[VN-MT-6, VN-22]%[0H4H2-1]
	Cho tam giác $ABC$, $R$ là bán kính đường tròn ngoại tiếp, chọn khẳng định đúng trong các khẳng định sau.
	\choice
	{$AB^2=AC^2+BC^2-AB\cdot BC\cdot\cos A$}
	{\True $\dfrac{AB}{\sin C}=\dfrac{BC}{\sin A}=\dfrac{AC}{\sin B}=2R$}
	{$AB^2=AC^2+BC^2-2\cdot AB\cdot BC\cdot\cos A$}
	{$\dfrac{AC}{\sin B}=R$}
	\loigiai{
		\begin{center}
			\begin{tikzpicture}[line join = round,line cap = round, font = \footnotesize, scale = 1]
				\path
				(0:0) coordinate (B)
				+(0:5) coordinate (C)
				+(70:3) coordinate (A)
				($(A)!.5!(B)$) coordinate (M)
				($(B)!.5!(C)$) coordinate (N)
				($(A)!.5!(C)$) coordinate (P)
				($(M)!1mm!-90:(A)$) coordinate (M1)
				($(N)!1mm!-90:(B)$) coordinate (N1)
				($(P)!1mm!-90:(C)$) coordinate (P1)
				(intersection of M--M1 and N--N1) coordinate (O)
				;
				\draw
				let \p1=($(O)-(A)$) in (O) circle ({veclen(\x1,\y1)})
				(A)--(B)--(C)--cycle
				(O)--(A) node[midway,below left] {$R$} (O)--(B) (O)--(C)
				;
				\foreach \x/\g in {B/180,C/0,A/90,O/60}
				\fill (\x) circle (1pt)
				+(\g:3mm) node{$\x$};
			\end{tikzpicture}
			
		\end{center}
		Theo định lý sin ta có $\dfrac{AB}{\sin C}=\dfrac{BC}{\sin A}=\dfrac{AC}{\sin B}=2R$.
	}
\end{ex}

%%==========Câu 4
\begin{ex}%[0-HK1-CD-3-DongAnh-HaNoi-2324]%[VN-MT-6, VN-22]%[0H5H3-2]
	Cho tam giác $ABC$ có trọng tâm $G$, $M$ là trung điểm $BC$, mệnh đề nào sau đây đúng?
	\choice
	{\True $\overrightarrow{GB}+\overrightarrow{GC}=2\overrightarrow{GM}$}
	{$\overrightarrow{GA}+\overrightarrow{GB}=\overrightarrow{GC}$}
	{$\overrightarrow{GB}+\overrightarrow{GC}=2\overrightarrow{GA}$}
	{$\overrightarrow{AB}+\overrightarrow{AC}=2\overrightarrow{AG}$}
	\loigiai{
		\begin{center}
			\begin{tikzpicture}[line join = round,line cap = round, font = \footnotesize, scale = 1]
				\path
				(0:0) coordinate (B)
				+(0:5) coordinate (C)
				+(70:3) coordinate (A)
				($(A)!.5!(B)$) coordinate (N)
				($(B)!.5!(C)$) coordinate (M)
				($(A)!.5!(C)$) coordinate (P)
				(intersection of A--M and B--P) coordinate (G)
				;
				\draw
				(A)--(B)--(C)--cycle
				(A)--(M)
				(B)--(G)--(C)
				;
				\foreach \x/\g in {B/180,C/0,A/90,M/-90,G/70}
				\fill (\x) circle (1pt)
				+(\g:3mm) node{$\x$};
			\end{tikzpicture}
		\end{center}
		Ta có $\overrightarrow{GB}+\overrightarrow{GC}=\overrightarrow{GM}+\overrightarrow{MB}+\overrightarrow{GM}+\overrightarrow{MC}=2\overrightarrow{GM}+\overrightarrow{MB}+\overrightarrow{MC}$.\\
		Vì $M$ là trung điểm của $BC$ nên ta có $\overrightarrow{MB}+\overrightarrow{MC}=\overrightarrow{0}$.\\
		Vậy $\overrightarrow{GB}+\overrightarrow{GC}=2\overrightarrow{GM}$.
	}
\end{ex}

%%==========Câu 5
\begin{ex}%[0-HK1-CD-3-DongAnh-HaNoi-2324]%[VN-MT-6, VN-22]%[0H5H3-2]
	Cho ba điểm phân biệt $A$, $B$, $C$. Nếu $\overrightarrow{AB}=-3\overrightarrow{AC}$ thì đẳng thức nào dưới đây đúng?\\
	\choice
	{$\overrightarrow{BC}=-2\overrightarrow{AC}$}
	{$\overrightarrow{BC}=2\overrightarrow{AC}$}
	{$\overrightarrow{BC}=-4\overrightarrow{AC}$}
	{\True $\overrightarrow{BC}=4\overrightarrow{AC}$}
	\loigiai{
		\begin{center}
			\begin{tikzpicture}[scale=1,font=\footnotesize,line join=round, line cap = round, >=stealth]
				\path
				(0:0) coordinate (A)
				+(0:3) coordinate (B)
				+(180:1) coordinate (C)
				;
				\draw[-stealth](A)--(B);
				\draw[-stealth](A)--(C);
				\foreach \x/\g in {B/0,C/180,A/90}
				\fill (\x) circle (1pt)
				+(\g:3mm) node{$\x$};
			\end{tikzpicture}
		\end{center}
		Từ hình vẽ ta thấy $\overrightarrow{BC}=4\overrightarrow{AC}$.
	}
\end{ex}

%%==========Câu 6
\begin{ex}%[0-HK1-CD-3-DongAnh-HaNoi-2324]%[VN-MT-6, VN-22]%[0D7H1-2]
	Cho $f(x)=x^2-4$. Tập hợp $S$ gồm các giá trị của $x$ để $f(x) < 0$ là
	\choice
	{$S=(-2;+\infty)$}
	{$S=\mathbb{R}\setminus\left[-2;2\right]$}
	{$S=(-\infty;-2)$}
	{\True $S=(-2;2)$}
	\loigiai{
		Ta có $f(x) < 0\Leftrightarrow x^2-4<0 \Leftrightarrow x^2<4 \Leftrightarrow -2<x<2$.\\
		Vậy tập hợp $S=(-2;2)$.
		
	}
\end{ex}

%%==========Câu 7
\begin{ex}%[0-HK1-CD-3-DongAnh-HaNoi-2324]%[VN-MT-6, VN-22]%[0D3H2-4]
	Tọa độ các giao điểm của parabol $(P)\colon y=x^2+5x+4$ với trục hoành là
	\choice
	{\True $(-1;0)$; $(-4;0)$}
	{$(0;-1)$; $(0;-4)$}
	{$(-1;0)$; $(0;-4)$}
	{$(0;-1)$; $(-4;0)$}
	\loigiai{
		Hoành độ giao điểm của $(P)$ với trục hoành là nghiệm của phương trình
		\[x^2+5x+4=0\Leftrightarrow \hoac{&x=-1\\&x=-4.}\]
		Vậy tọa độ giao điểm của $(P)$ với trục hoành là $(-1;0)$; $(-4;0)$.
	}
\end{ex}

%%==========Câu 8
\begin{ex}%[0-HK1-CD-3-DongAnh-HaNoi-2324]%[VN-MT-6, VN-22]%[0H5H2-1]
	Cho hình bình hành $ABCD$. Chọn khẳng định đúng.
	\choice
	{$\overrightarrow{BA}=\overrightarrow{DC}$}
	{$\overrightarrow{CA}=\overrightarrow{CB}-\overrightarrow{BA}$}
	{\True $\overrightarrow{BA}+\overrightarrow{BC}=\overrightarrow{BD}$}
	{$\overrightarrow{AC}=\overrightarrow{BA}+\overrightarrow{BC}$}
	\loigiai{
		\begin{center}
			\begin{tikzpicture}[scale=1,font=\footnotesize,line join=round, line cap = round, >=stealth]
				\path
				(0:0) coordinate (D)
				+(0:5) coordinate (C)
				+(65:3) coordinate (A)
				($(A)+(C)-(D)$) coordinate (B)
				;
				\draw
				(A)--(B)--(C)--(D)--cycle
				;
				\foreach \x/\g in {D/-90,C/-90,A/90,B/90}
				\fill (\x) circle (1pt)
				+(\g:3mm) node{$\x$};
			\end{tikzpicture}
		\end{center}
		Theo quy tắc hình bình hành
		$\overrightarrow{BA}+\overrightarrow{BC}=\overrightarrow{BD}$.
	}
\end{ex}

%%==========Câu 9
\begin{ex}%[0-HK1-CD-3-DongAnh-HaNoi-2324]%[VN-MT-6, VN-22]%[0H5H2-4]
	Cho hình vuông $ABCD$ có cạnh bằng $a$. Độ dài $\left|\overrightarrow{AD}+\overrightarrow{AB}\right|$ bằng
	\choice
	{$\dfrac{a\sqrt{3}}{2}$}
	{\True $a\sqrt{2}$}
	{$2a$}
	{$\dfrac{a\sqrt{2}}{2}$}
	\loigiai{
		\begin{center}
			\begin{tikzpicture}[scale=1,font=\footnotesize,line join=round, line cap = round, >=stealth]
				\def\canh{3}
				\path
				(0:0) coordinate (A)
				++(0:\canh) coordinate (B)
				++(-90:\canh) coordinate (C)
				($(A)+(C)-(B)$) coordinate (D)
				(A)--(B) node[midway,above] {$a$}
				;
				\draw (A)--(B)--(C)--(D)--cycle
				;
				\foreach \x/\g in {A/90,B/90,C/-90,D/-90}
				\fill (\x) circle (1pt)
				+(\g:3mm) node {$\x$};
			\end{tikzpicture}
		\end{center}
		Ta có $\overrightarrow{AD}+\overrightarrow{AB}=\overrightarrow{AC} \Rightarrow\left|\overrightarrow{AD}+\overrightarrow{AB}\right|=\left|\overrightarrow{AC}\right|=AC$.\\
		Mà $AC=a\sqrt{2} \Rightarrow \left|\overrightarrow{AD}+\overrightarrow{AB}\right|=a\sqrt{2}$.
	}
\end{ex}

%%==========Câu 10
\begin{ex}%[0-HK1-CD-3-DongAnh-HaNoi-2324]%[VN-MT-6, VN-22]%[0D7H3-1]
	Tập nghiệm của phương trình $\sqrt{x^2+5}=\sqrt{2x^2+1}$ là
	\choice
	{$S=\left\{0;-2\right\}$}
	{$S=\{-2\}$}
	{\True $S=\left\{-2;2\right\}$}
	{$S=\{0\}$}
	\loigiai{
		Ta có $x^2+5>0$.\\
		Phương trình tương đương với
		\begin{align*}
			&x^2+5=2x^2+1\\
			\Leftrightarrow~&x^2=4\\
			\Leftrightarrow~&x=\pm2.
		\end{align*}
		Vậy tập nghiệm của phương trình là $S=\left\{-2;2\right\}$.
		
	}
\end{ex}

%%==========Câu 11
\begin{ex}%[0-HK1-CD-3-DongAnh-HaNoi-2324]%[VN-MT-6, VN-22]%[0D3H2-2]
	Cho $f(x)=a{x^2}+bx+c~(a\ne 0)$ có bảng xét dấu dưới đây.
	\begin{center}
		\begin{tikzpicture}
			\tkzTabInit[nocadre=true, lgt=1.2, espcl=3, deltacl=0.6]
			{$x$/0.8,$f(x)$/0.6}
			{$-\infty$,$x_1$,$x_2$,$+\infty$};
			\tkzTabLine{,+,$0$,-,$0$,+,};
			\path
			($(N11)!.5!(N21)$) coordinate (M) node[above=1.5mm] {$0$}
			;
		\end{tikzpicture}
	\end{center}
	Hỏi mệnh đề nào dưới đây đúng?
	\choice
	{$a>0$, $b>0$, $c>0$}
	{\True $a>0$, $b<0$, $c>0$}
	{$a>0$, $b<0$, $c<0$}
	{$a<0$, $b<0$, $c>0$}
	\loigiai{
		Từ bảng xét dấu ta thấy đồ thị là dạng Parabol có bề lõm quay lên trên nên $a>0$.\\
		Hoành độ đỉnh của Parabol dương nên $-\dfrac{b}{2a}>0 \Rightarrow b<0$.\\
		Parabol cắt trục tung tại điểm tung độ dương nên $f(0)>0\Leftrightarrow c>0$.
	}
\end{ex}

%%==========Câu 12
\begin{ex}%[0-HK1-CD-3-DongAnh-HaNoi-2324]%[VN-MT-6, VN-22]%[0D7H1-2]
	Cho tam thức bậc hai $f(x)=2x^2+x-1$. Tìm $x$ để $f(x) > 0$.
	\choice
	{$x\in\left(-\infty;-1\right]\cup\left[\dfrac{1}{2};+\infty\right)$}
	{$x\in\left(-\infty;-1\right)\cap\left(\dfrac{1}{2};+\infty\right)$}
	{\True $x\in\left(-\infty;-1\right)\cup\left(\dfrac{1}{2};+\infty\right)$}
	{$x\in\left(-1;\dfrac{1}{2}\right)$}
	\loigiai{
		Xét $f(x)=2x^2+x-1=0$ có hai nghiệm $x=-1$; $x=\dfrac{1}{2}$.\\
		Để $f(x)>0$	thì $x<-1$ hoặc $x>\dfrac{1}{2}$.\\
		Vậy $x\in (-\infty;-1)\cup \left(\dfrac{1}{2};+\infty\right)$.
		
	}
\end{ex}

\Closesolutionfile{ans}

%================================PHẦN II
\cauds
\Opensolutionfile{ans}[ans/0-HK1-CD-3-DongAnh-HaNoi-2324-Phan-II]

%%==========Câu 13
\begin{ex}%[0-HK1-CD-3-DongAnh-HaNoi-2324]%[VN-MT-6, VN-22]%[0H5V3-5]
	Cho tam giác $ABC$ có trọng tâm $G$. Điểm $M$ là trung điểm của $BC$.
	\choiceTF
	{\True $\overrightarrow{MB}+\overrightarrow{MC}=\overrightarrow{0}$}
	{\True $\overrightarrow{GA}+\overrightarrow{GB}+\overrightarrow{GC}=\overrightarrow{0}$}
	{$\overrightarrow{MA}+\overrightarrow{MB}+\overrightarrow{MC}=\overrightarrow{0}$}
	{\True $\overrightarrow{AG}=\dfrac{1}{3}\left(\overrightarrow{AB}+\overrightarrow{AC}\right)$}
	\loigiai{
		\begin{center}
			\begin{center}
				\begin{tikzpicture}[scale=1,font=\footnotesize,line join=round, line cap = round, >=stealth]
					\path
					(0:0) coordinate (B)
					+(0:5) coordinate (C)
					+(70:3) coordinate (A)
					($(A)!.5!(B)$) coordinate (N)
					($(B)!.5!(C)$) coordinate (M)
					($(A)!.5!(C)$) coordinate (P)
					(intersection of A--M and B--P) coordinate (G)
					;
					\draw
					(A)--(B)--(C)--cycle
					(A)--(M)
					(B)--(G)--(C)
					;
					\foreach \x/\g in {B/180,C/0,A/90,M/-90,G/70}
					\fill (\x) circle (1pt)
					+(\g:3mm) node{$\x$};
				\end{tikzpicture}
			\end{center}
		\end{center}
		\begin{itemchoice}
			\itemch \textbf{Đúng.}\\
			$M$ là trung điểm của $BC$ nên $\overrightarrow{MB}+\overrightarrow{MC}=\overrightarrow{0}$.
			\itemch \textbf{Đúng.}\\
			$G$ là trọng tâm của tam giác $ABC$ nên $\overrightarrow{GA}+\overrightarrow{GB}+\overrightarrow{GC}=\overrightarrow{0}$.
			\itemch \textbf{Sai.}\\ $\overrightarrow{MA}+\overrightarrow{MB}+\overrightarrow{MC}=\overrightarrow{MA}$ (vì $\overrightarrow{MB}+\overrightarrow{MC}=\overrightarrow{0}$).
			\itemch \textbf{Đúng.}\\
			$\overrightarrow{AG}=\dfrac{2}{3}\overrightarrow{AM}=\dfrac{1}{3}\left(\overrightarrow{AB}+\overrightarrow{AC}\right)$.
		\end{itemchoice}
		
		
	}
\end{ex}

%%==========Câu 14
\begin{ex}%[0-HK1-CD-3-DongAnh-HaNoi-2324]%[VN-MT-6, VN-22]%[0D3V1-4]
	Cho đồ thị hàm số $y=ax^2+bx+c$.
	\begin{center}
		\begin{tikzpicture}[scale=1,font=\footnotesize,line join=round, line cap = round, >=stealth]
			%Gán số liệu.
			\def\xmin{-4};\def\ymin{-4};\def\xmax{4};\def\ymax{4};
			%Gán tọa độ.
			\coordinate (O) at (0,0);
			%Trục Oxy.
			\draw[->] (\xmin,0)--(\xmax,0) node[below]{$x$};
			\draw[->] (0,\ymin)--(0,\ymax) node[left]{$y$};
			\fill (O) node[below left]{$O$} circle(1pt);
			%Giới hạn đồ thị.
			\clip ({\xmin-0.1},{\ymin-0.1}) rectangle ({\xmax+0.1},{\ymax+0.1});
			\foreach \x in {-3,-2,-1,2,3}{
				\fill (\x,0) node[below]{$\x$} circle(1pt);
			}
			\foreach \x in {1}{
				\fill (\x,0) node[below right]{$\x$} circle(1pt);
			}
			\foreach \y in {-3,-2,-1,1,2,3}{
				\fill (0,\y) node[left]{$\y$} circle(1pt);
			}
			\fill (1,-3) circle (1pt);
			\draw[dashed] (0,-3)-|(1,0);
			%Vẽ đồ thị.
			\draw[samples=100] plot[domain=-4:4](\x,{(\x)^2-2*\x-2});
		\end{tikzpicture}
	\end{center}
	\choiceTF
	{\True Đồ thị hàm số trên có dạng Parabol}
	{Đồ thị có một tâm đối xứng tọa độ $(1;-3)$}
	{\True Đồ thị trên là đồ thị của hàm số $y=x^2-2x-2$}
	{Phương trình $|x^2-2x-2|=m$ có $4$ nghiệm phân biệt khi $0\leq m \leq 3$}
	\loigiai{
		\begin{itemchoice}
			\itemch \textbf{Đúng.}\\
			Đồ thị hàm số trên có dạng Parabol.
			\itemch \textbf{Sai.}\\
			Đồ thị không có tâm đối xứng.
			\itemch \textbf{Đúng.}\\
			Giả sử đồ thị trên của hàm số $y=ax^2+bx+c$.\\
			Đồ thị cắt trục tung tại điểm có tọa độ $(0;-2)$ nên $c=-2$.\\
			Đồ thị có đỉnh tại $x=1\Rightarrow-\dfrac{b}{2a}=1\Rightarrow b=-2a.\qquad (1)$\\
			Điểm $(1;-3)$ thuộc đồ thị hàm số nên ta có $-3=a+b-2 \Rightarrow a+b=-1.\qquad (2)$\\
			Từ $(1)$ và $(2)$ suy ra $a=1$; $b=-2$.\\
			Vậy đồ thị trên là của hàm số $y=x^2-2x-2$.
			\itemch \textbf{Sai.}\\
			Hàm số 
			$y=\heva{&{x^2-2x-2~\text{với}~x \in \left(-\infty;1-\sqrt{2}\right] \cup \left[1+\sqrt{2};+\infty\right)}\\&-x^2+2x+2~\text{với}~x\in \left(1-\sqrt{2};1+\sqrt{2}\right).}$\\
			Ta có đồ thị hàm số
			\begin{center}
				\begin{tikzpicture}[scale=1,font=\footnotesize,line join=round, line cap = round, >=stealth]
					%Gán số liệu.
					\def\xmin{-3};\def\ymin{-1};\def\xmax{4};\def\ymax{4};
					%Gán tọa độ.
					\coordinate (O) at (0,0);
					%Trục Oxy.
					\draw[->] (\xmin,0)--(\xmax,0) node[below]{$x$};
					\draw[->] (0,\ymin)--(0,\ymax) node[left]{$y$};
					\fill (O) node[below left]{$O$} circle(1pt);
					%Giới hạn đồ thị.
					\clip ({\xmin-0.1},{\ymin-0.1}) rectangle ({\xmax+0.1},{\ymax+0.1});
					\foreach \x in {-2,-1,1,2,3}{
						\fill (\x,0) node[below]{$\x$} circle(1pt);
					}
					\foreach \y in {1,2,3}{
						\fill (0,\y) node[left]{$\y$} circle(1pt);
					}
					\fill (1,3) circle (1pt);
					\draw[dashed] (0,3)-|(1,0);
					%Vẽ đồ thị.
					\draw[samples=100] plot[domain=-3:4](\x,{abs((\x)^2-2*(\x)-2)});
				\end{tikzpicture}
			\end{center}
			Số nghiệm của phương trình $|x^2-2x-2|=m~(*)$ chính là số giao điểm của hai đồ thị hàm số $y=|x^2-2x-2|$ và $y=m$.\\
			Từ đồ thị ta thấy phương trình $(*)$ có $4$ nghiệm phân biệt khi và chỉ khi $0<m<3$.
		\end{itemchoice}		
	}
\end{ex}

%%==========Câu 15
\begin{ex}%[0-HK1-CD-3-DongAnh-HaNoi-2324]%[VN-MT-6, VN-22]%[0D7V1-2]
	Cho hàm số $f(x)=x^2+2mx+3m-2$.
	\choiceTF
	{\True Tập xác định của hàm số $f(x)$ là $\mathscr{D}=\mathbb{R}$}
	{$f(x)$ là tam thức bậc hai có $\Delta'=m^2+3m-2$}
	{$f(x)\geq 0 \Leftrightarrow \heva{&a>0\\&\Delta'<0}$}
	{\True Với $m\in\left[1;2\right]$ thì $f(x)\geq 0,\forall x\in \mathbb{R}$}
	\loigiai{
		\begin{itemchoice}
			\itemch \textbf{Đúng.} \\
			Hàm $f(x)$ là hàm đa thức nên tập xác định của hàm số $\mathscr{D}=\mathbb{R}$.
			\itemch \textbf{Sai.} \\
			$f(x)$ là tam thức bậc hai có $\Delta'=m^2-3m+2$.
			\itemch \textbf{Sai.} \\
			Ta có $f(x)\geq 0 \Leftrightarrow \heva{&a>0\\&\Delta'\leq 0.}$
			\itemch \textbf{Đúng.} \\
			Ta có $\Delta'\leq 0 \Leftrightarrow m^2-3m+2 \leq 0 \Leftrightarrow 1\leq m\leq 2$.\\
			Vậy với $m\in\left[1;2\right]$ thì $f(x)\geq 0$, $\forall x\in \mathbb{R}$.
		\end{itemchoice}
		
		
	}
\end{ex}

%%==========Câu 16
\begin{ex}%[0-HK1-CD-3-DongAnh-HaNoi-2324]%[VN-MT-6, VN-22]%[0D7V3-4]
	Cho phương trình $\sqrt{x^2-x+1}+\sqrt{x^2+x+1}=\sqrt{4-x}$. Gọi $S$ là tập hợp các nghiệm của phương trình.
	\choiceTF
	{Điều kiện $x\geq 4$}
	{\True Phương trình đã cho tương đương với $\sqrt{x^2-x+1}=\sqrt{4-x}-\sqrt{x^2+x+1}$}
	{Phương trình có $3$ nghiệm phân biệt}
	{Tổng các phần tử của $S$ là $\dfrac{11-\sqrt{185}}{10}$}
	\loigiai{
		\begin{itemchoice}
			\itemch \textbf{Sai.}\\
			Điều kiện $x\leq 4$.
			\itemch \textbf{Đúng.}\\
			Phương trình tương đương với $\sqrt{x^2-x+1}=\sqrt{4-x}-\sqrt{x^2+x+1}$.
			\itemch \textbf{Sai.} \\
			Bình phương hai vế của phương trình ta có
			\begin{align*}
				4+x=2 \sqrt{(4-x)\left(x^2+x+1\right)} \Leftrightarrow 
				4 x^3-11 x^2-4 x=0 \Leftrightarrow \hoac{&x=0\\&x=\dfrac{11-\sqrt{185}}{8}\\&x=\dfrac{11+\sqrt{185}}{8}}.
			\end{align*}
			Thay các nghiệm vào phương trình ta thấy có hai nghiệm $x=0$ và $x=\dfrac{11-\sqrt{185}}{8}$ thỏa mãn.
			Vậy phương trình có hai nghiệm phân biệt.
			\itemch \textbf{Sai.} \\
			Tập hợp các nghiệm của phương trình là $S=\left\{0;\dfrac{11-\sqrt{185}}{8}\right\}$.\\
			Tổng các phần tử của $S$ là $\dfrac{11-\sqrt{185}}{8}$.
		\end{itemchoice}
		
		
	}
\end{ex}

\Closesolutionfile{ans}

%================================PHẦN III
\caukq
\Opensolutionfile{ans}[ans/0-HK1-CD-3-DongAnh-HaNoi-2324-Phan-III]

\begin{ex}%[0-HK1-CD-3-DongAnh-HaNoi-2324]%[VN-MT-6, VN-22]%[0D7H2-3]
	Bất phương trình $\left(x^2-4x+4\right)\left(-x^2-2x+3\right)\ge 0$ có tập nghiệm $S=[a;b]\cup\{2\}$. Tính tổng $a+b$.	
	\shortans[]{$-2$}
	\loigiai{
		Ta có
		$(x^2-4x+4)(-x^2-2x+3)\geq 0 
		\Leftrightarrow (x-2)^2(-x^2-2x+3) \geq 0$.\\
		Trường hợp 1: $x=2\Leftrightarrow x-2=0$ bất phương trình nghiệm đúng.\\
		Trường hợp 2: $x\ne2\Leftrightarrow (x-2)^2>0$ ta có $-x^2-2x+3\geq 0 \Leftrightarrow -3\leq x\leq 1$.\\
		Vậy tập nghiệm của bất phương trình là $S=[-3;1]\cup \{2\}$.\\
		Vậy $a=-3$; $b=1$ nên $a+b=-2$.
	}
\end{ex}

\begin{ex}%[0-HK1-CD-3-DongAnh-HaNoi-2324]%[VN-MT-6, VN-22]%[0H5H4-1]
	Cho tam giác $ABC$ có $\widehat{A}=60^\circ$, $AB=5$, $AC=8$. Tính $\overrightarrow{BC}\cdot\overrightarrow{AC}$.
	\shortans[]{$44$}
	\loigiai{
		Ta có \[\overrightarrow{BC}\cdot\overrightarrow{AC}=\big(\overrightarrow{AC}-\overrightarrow{AB}\big)\cdot \overrightarrow{AC} =\overrightarrow{AC}^2-\overrightarrow{AB}\cdot\overrightarrow{AC}=8^2-5\cdot 8\cdot \cos 60^\circ=44.\]		
	}
\end{ex}
\begin{ex}%[0-HK1-CD-3-DongAnh-HaNoi-2324]%[VN-MT-6, VN-22]%[0D3H1-4]
	Hoành độ đỉnh của parabol $y=x^2-4x+3$ bằng bao nhiêu?
	\shortans[]{$2$}
	\loigiai{
		Đồ thị hàm số $y=x^2-4x+3$ là parabol có hoành độ đỉnh $x=-\dfrac{b}{2a}=\dfrac{4}{2}=2$.		
	}
\end{ex}
\begin{ex}%[0-HK1-CD-3-DongAnh-HaNoi-2324]%[VN-MT-6, VN-22]%[0H5H4-3]
	Cho hình vuông $ABCD$ tâm $O$ có cạnh bằng $1$. Tính $\overrightarrow{AB}\cdot\overrightarrow{AC}+\overrightarrow{DO}\cdot\overrightarrow{CA}+\overrightarrow{AB}\cdot\overrightarrow{BO}$.
	
	\shortans[]{$0{,}5$}
	\loigiai{
		\begin{center}
			\begin{tikzpicture}[scale=1,font=\footnotesize,line join=round, line cap = round, >=stealth]
				%Gán số liệu.
				\def\canhAB{3};
				%Gán tọa độ.
				\coordinate (B) at (0,0);
				\coordinate (A) at ($(B)+(90:\canhAB)$);
				\coordinate (C) at ($(B)+(0:\canhAB)$);
				\coordinate (D) at ($(A)+(0:\canhAB)$);
				\path (intersection of A--C and B--D) coordinate (O);				
				%Vẽ hình vuông ABCD.
				\draw (A) rectangle (C)
				(A)--(C) (B)--(D)
				;
				%Gán nhãn.
				\foreach \x/\y in {A/90,B/-90,C/-90,D/90,O/90}{\fill (\x) circle(1pt) ($(\x)+(\y:0.3cm)$) node{$\x$};}
			\end{tikzpicture}
		\end{center}
		Vì $DO\perp CA$ nên $\overrightarrow{DO}\cdot \overrightarrow{CA}=0$.\\
		$\overrightarrow{AB}\cdot\overrightarrow{AC}=AB\cdot AC \cdot \cos 45^\circ=1\cdot \sqrt{2}\cdot \dfrac{\sqrt{2}}{2}=1$.\\
		$\overrightarrow{AB}\cdot\overrightarrow{BO}=AB\cdot BO \cdot \cos 135^\circ=1\cdot \dfrac{\sqrt{2}}{2}\cdot \left(-\dfrac{\sqrt{2}}{2}\right)=-\dfrac{1}{2}$.\\
		Vậy $\overrightarrow{AB}\cdot\overrightarrow{AC}+\overrightarrow{DO}\cdot\overrightarrow{CA}+\overrightarrow{AB}\cdot\overrightarrow{BO}=1-\dfrac{1}{2}=\dfrac{1}{2}=0{,}5$.
	}
\end{ex}

\begin{ex}%[0-HK1-CD-3-DongAnh-HaNoi-2324]%[VN-MT-6, VN-22]%[0D7H3-5]
	Phương trình $x^4-24x^2-25=0$ có bao nhiêu nghiệm thực?
	\shortans[]{$2$}
	\loigiai{
		Ta xét $x^4-24x^2-25=0$.\\
		Đặt $t=x^2$, ($t>0$) ta có $t^2-24t-25=0.\qquad (1)$\\
		Phương trình $(1)$ có hai nghiệm $t=-1$ (loại); $t=25$.\\
		Vậy phương trình có hai nghiệm thực $x=\pm 5$.
	}
\end{ex}

\begin{ex}%[0-HK1-CD-3-DongAnh-HaNoi-2324]%[VN-MT-6, VN-22]%[0D7C3-4]
	Cho phương trình $\sqrt{\dfrac{x^3+1}{x+3}}+\sqrt{x+3}=\sqrt{x^2-x+1}+\sqrt{x+1}$. Gọi $S$ là tập nghiệm của phương trình, tổng các phần tử của $S$ bằng bao nhiêu?
	\shortans[]{$2$}
	\loigiai{
		Điều kiện $x\ge-1$. \\
		Bình phương hai vế của phương trình ta có
		\begin{align*}
			&\left(\sqrt{\dfrac{x^3+1}{x+3}}+\sqrt{x+3}\right)^2=\left(\sqrt{x^2-x+1}+\sqrt{x+1}\right)^2 \\\Leftrightarrow~&\dfrac{x^3+1}{x+3}+2 \sqrt{x^3+1}+(x+3)=\left(x^2-x+1\right)+2 \sqrt{(x+1)\left(x^2-x+1\right)}+(x+1)\\ \Leftrightarrow~&\dfrac{x^3+1}{x+3}=x^2-x-1\\ \Leftrightarrow~&x^2-2 x-2=0\\ \Leftrightarrow~& x=1 \pm\sqrt{3}.
		\end{align*}
		Thử lại ta thấy $x=1\pm \sqrt{3}$ thỏa mãn phương trình đã cho nên phương trình có tập nghiệm là $S=\left\{1-\sqrt{3};1+\sqrt{3}\right\}$.\\
		Tổng các phần tử của tập hợp nghiệm $S$ bằng $2$.
	}
\end{ex}

\Closesolutionfile{ans}
% \begin{indapan}
% 	{ans/0-HK1-CD-3-DongAnh-HaNoi-2324}
% \end{indapan}


